% GENERATED FILE. Edit canonical lessons, then run npm run generate:books.
% canonical-source-hash: fnv1a64:99a4ad6c679d2518
% canonical-lessons: LA-C503-natura, LA-C503-indoles, LA-C503-alacer, LA-C503-lepidus, LA-C503-lascivus

\chapter{Temperament}
\label{ch:la-temperament}
\hlchaptermodality{503}

\begin{chapteropening}
\textbf{By the end of this chapter:} \emph{I can describe someone's nature and temperament.}

Complete the last lesson of chapter 503: I can describe someone's nature and temperament.
\end{chapteropening}

\section[nātūra, nātūrae]{nātūra, nātūrae — nature, character}

\label{lesson:LA-C503-natura}



\begin{quote}
Before the new one: say the Latin for funny, then the Latin for witty.
\end{quote}

\subsection*{You'll want to know: nātūra, nātūrae}
\textbf{nātūra, nātūrae} — \textquotedblleft{}nature, character\textquotedblright{}.

\begin{quote}
Nātūra frātris meī placida est.
\end{quote}

\begin{quote}
\emph{My brother's nature is calm.}
\end{quote}

Feminine, first declension, from \emph{nāscor} — to be born: what you are born with. English \emph{nature}.

\begin{cousinweb}
\textbf{nātūra} is built on \emph{nātus}, the participle of \emph{nāscor} \textquotedblleft{}to be born\textquotedblright{} (Lewis \& Short; EtymDB). English \emph{nature} comes from it.
\end{cousinweb}

\begin{grammarlens}[title={What you've built}]
One more word for telling, explaining and describing at B1.
\end{grammarlens}

\subsection*{Your turn}
\begin{itemize}
\raggedright
  \item \emph{Say it:} \emph{nātūra, nātūrae}
  \item \emph{Say it:} \emph{nātūra, nātūrae}, once more
  \item \emph{Say it:} a sentence of your own with \emph{nātūra, nātūrae}
  \item \emph{Recall:} say the Latin for funny, then the Latin for witty, then say \emph{nātūra, nātūrae} again
\end{itemize}

\begin{tcolorbox}[breakable,colback=teal!4,colframe=teal!35!black,title={Before you move on}]
What does nātūra, nātūrae mean? (\textquotedblleft{}Nature, character\textquotedblright{}.) Say it once more.
\end{tcolorbox}
\section[indolēs, indolis]{indolēs, indolis — natural disposition, talent}

\label{lesson:LA-C503-indoles}



\begin{quote}
Before the new one: say the Latin for witty, then the Latin for nature.
\end{quote}

\subsection*{You'll want to know: indolēs, indolis}
\textbf{indolēs, indolis} — \textquotedblleft{}natural disposition, talent\textquotedblright{}.

\begin{quote}
Hic puer bonam indolem habet.
\end{quote}

\begin{quote}
\emph{This boy has a good disposition.}
\end{quote}

Feminine, third declension. Inborn character, often of a promising young person.

\begin{cousinweb}
\textbf{indolēs} is built on old \emph{indu} \textquotedblleft{}within\textquotedblright{} and the root of \emph{olēscō} \textquotedblleft{}to grow\textquotedblright{}, a relative of \emph{alō} \textquotedblleft{}to nourish\textquotedblright{} (Lewis \& Short; EtymDB): what grows from within.
\end{cousinweb}

\begin{grammarlens}[title={What you've built}]
One more word for telling, explaining and describing at B1.
\end{grammarlens}

\subsection*{Your turn}
\begin{itemize}
\raggedright
  \item \emph{Say it:} \emph{indolēs, indolis}
  \item \emph{Say it:} \emph{indolēs, indolis}, once more
  \item \emph{Say it:} a sentence of your own with \emph{indolēs, indolis}
  \item \emph{Recall:} say the Latin for witty, then the Latin for nature, then say \emph{indolēs, indolis} again
\end{itemize}

\begin{tcolorbox}[breakable,colback=teal!4,colframe=teal!35!black,title={Before you move on}]
What does indolēs, indolis mean? (\textquotedblleft{}Natural disposition, talent\textquotedblright{}.) Say it once more.
\end{tcolorbox}
\section[alacer, alacris, alacre]{alacer, alacris, alacre — lively, eager}

\label{lesson:LA-C503-alacer}



\begin{quote}
Before the new one: say the Latin for nature, then the Latin for natural disposition.
\end{quote}

\subsection*{You'll want to know: alacer, alacris, alacre}
\textbf{alacer, alacris, alacre} — \textquotedblleft{}lively, eager\textquotedblright{}.

\begin{quote}
Mīlitēs alacrēs ad pugnam currunt.
\end{quote}

\begin{quote}
\emph{The soldiers run eagerly to the fight.}
\end{quote}

Three endings in the nominative, like \emph{ācer}. English \emph{alacrity} comes from \emph{alacritās}.

\begin{cousinweb}
\textbf{alacer} is a native Latin adjective; Lewis \& Short suggests a link with \emph{alō} \textquotedblleft{}to nourish\textquotedblright{}, but the origin is not settled. English \emph{alacrity} comes from \emph{alacritās}, built on it.
\end{cousinweb}

\begin{grammarlens}[title={What you've built}]
One more word for telling, explaining and describing at B1.
\end{grammarlens}

\subsection*{Your turn}
\begin{itemize}
\raggedright
  \item \emph{Say it:} \emph{alacer, alacris, alacre}
  \item \emph{Say it:} \emph{alacer, alacris, alacre}, once more
  \item \emph{Say it:} a sentence of your own with \emph{alacer, alacris, alacre}
  \item \emph{Recall:} say the Latin for nature, then the Latin for natural disposition, then say \emph{alacer, alacris, alacre} again
\end{itemize}

\begin{tcolorbox}[breakable,colback=teal!4,colframe=teal!35!black,title={Before you move on}]
What does alacer, alacris, alacre mean? (\textquotedblleft{}Lively, eager\textquotedblright{}.) Say it once more.
\end{tcolorbox}
\section[lepidus, lepida, lepidum]{lepidus, lepida, lepidum — charming, delightful}

\label{lesson:LA-C503-lepidus}



\begin{quote}
Before the new one: say the Latin for natural disposition, then the Latin for lively.
\end{quote}

\subsection*{You'll want to know: lepidus, lepida, lepidum}
\textbf{lepidus, lepida, lepidum} — \textquotedblleft{}charming, delightful\textquotedblright{}.

\begin{quote}
Fābula lepida et brevis est.
\end{quote}

\begin{quote}
\emph{The story is charming and short.}
\end{quote}

A word of comedy and of friendly letters: neat, pleasant, nicely done.

\begin{cousinweb}
\textbf{lepidus} is built on \emph{lepōs} \textquotedblleft{}charm, grace\textquotedblright{} (Lewis \& Short); the origin of \emph{lepōs} is not settled.
\end{cousinweb}

\begin{grammarlens}[title={What you've built}]
One more word for telling, explaining and describing at B1.
\end{grammarlens}

\subsection*{Your turn}
\begin{itemize}
\raggedright
  \item \emph{Say it:} \emph{lepidus, lepida, lepidum}
  \item \emph{Say it:} \emph{lepidus, lepida, lepidum}, once more
  \item \emph{Say it:} a sentence of your own with \emph{lepidus, lepida, lepidum}
  \item \emph{Recall:} say the Latin for natural disposition, then the Latin for lively, then say \emph{lepidus, lepida, lepidum} again
\end{itemize}

\begin{tcolorbox}[breakable,colback=teal!4,colframe=teal!35!black,title={Before you move on}]
What does lepidus, lepida, lepidum mean? (\textquotedblleft{}Charming, delightful\textquotedblright{}.) Say it once more.
\end{tcolorbox}
\section[lascīvus, lascīva, lascīvum]{lascīvus, lascīva, lascīvum — playful, frisky}

\label{lesson:LA-C503-lascivus}



\begin{quote}
Before the new one: say the Latin for lively, then the Latin for charming.
\end{quote}

\subsection*{You'll want to know: lascīvus, lascīva, lascīvum}
\textbf{lascīvus, lascīva, lascīvum} — \textquotedblleft{}playful, frisky\textquotedblright{}.

\begin{quote}
Catulī lascīvī in hortō lūdunt.
\end{quote}

\begin{quote}
\emph{The playful puppies play in the garden.}
\end{quote}

Of animals and children, frisky; of grown-ups, sometimes naughty. English \emph{lascivious} is narrower.

\begin{cousinweb}
\textbf{lascīvus} is inherited and traced to a Proto-Indo-European root \emph{las-} \textquotedblleft{}eager, wanton\textquotedblright{} (EtymDB; Lewis \& Short). English \emph{lascivious} comes from it, with a narrower meaning.
\end{cousinweb}

\begin{grammarlens}[title={What you've built}]
One more word for telling, explaining and describing at B1.
\end{grammarlens}

\subsection*{Your turn}
\begin{itemize}
\raggedright
  \item \emph{Say it:} \emph{lascīvus, lascīva, lascīvum}
  \item \emph{Say it:} \emph{lascīvus, lascīva, lascīvum}, once more
  \item \emph{Say it:} a sentence of your own with \emph{lascīvus, lascīva, lascīvum}
  \item \emph{Recall:} say the Latin for lively, then the Latin for charming, then say \emph{lascīvus, lascīva, lascīvum} again
\end{itemize}

\begin{tcolorbox}[breakable,colback=teal!4,colframe=teal!35!black,title={Before you move on}]
What does lascīvus, lascīva, lascīvum mean? (\textquotedblleft{}Playful, frisky\textquotedblright{}.) Say it once more.
\end{tcolorbox}
